\section{GlyphChain release-evidence artefact contract}
\label{sec:glyphchain-release-evidence-artifact-contract}

Production readiness is not proven by attaching an arbitrary file to a signed
approval.  Each of the thirty-four release gates now declares the exact classes of
artefact required to support its claim.  The evaluator verifies the task-specific
contract before accepting authority signatures, so a validly signed but irrelevant
report cannot close a gate.

\subsection{Artefact commitment format}

Every retained artefact is represented inside the signed evidence body by

\begin{equation}
  A_i = (\mathrm{id}_i,\mathrm{kind}_i,
         \operatorname{SHA256}(\mathrm{bytes}_i)).
\end{equation}

The identifier is a stable opaque reference used by the controlled evidence store;
it is not a local path and must not contain private document content.  Identifiers
are unique within an evidence entry.  Kinds use the fixed lowercase protocol
vocabulary below.  The digest is exactly sixty-four lowercase hexadecimal
characters and binds the retained bytes rather than a filename or description.

An entry accepts between one and thirty-two artefact commitments.  For task $t$,
let $K_t$ be the required kind set and $K_E$ the kinds present in the signed body.
The artefact gate is

\begin{equation}
  G_{A,t} = (K_t \subseteq K_E)
  \land \operatorname{unique}(\mathrm{id}_{1\ldots n})
  \land \bigwedge_{i=1}^{n}\operatorname{valid}(A_i).
\end{equation}

Additional supporting artefacts are permitted within the bounded list, but they do
not replace a missing required kind.  Empty identifiers, duplicate identifiers,
unknown syntax, missing digests and task-inappropriate evidence fail closed.

\subsection{Live delivery artefacts}

\begin{center}
\begin{tabular}{p{0.13\linewidth}p{0.75\linewidth}}
\hline
\textbf{Gate} & \textbf{Required artefact kinds} \\
\hline
\texttt{LD-01} & \texttt{resend\_callback\_configuration};
\texttt{resend\_authenticated\_delivery\_event} \\
\texttt{LD-02} & \texttt{microsoft\_scope\_grant};
\texttt{exchange\_trace\_access\_grant} \\
\texttt{LD-03} & \texttt{microsoft\_draft\_send\_trace} \\
\texttt{LD-04} & \texttt{delivery\_failure\_retry\_matrix} \\
\texttt{LD-05} & \texttt{production\_delivery\_qualification\_report} \\
\hline
\end{tabular}
\end{center}

The callback event must be authenticated by the provider adapter.  Microsoft
evidence must preserve the exact Internet Message-ID, recipient and trace outcome.
The retry matrix must cover the required delay, failure, rejection and newly
authorised retry cases rather than merely show successful submission.

\subsection{Independent trust-publication artefacts}

\begin{center}
\begin{tabular}{p{0.13\linewidth}p{0.75\linewidth}}
\hline
\textbf{Gate} & \textbf{Required artefact kinds} \\
\hline
\texttt{TR-01} & \texttt{trust\_release\_manifest};
\texttt{trust\_release\_publication\_record} \\
\texttt{TR-02} & \texttt{independent\_checkpoint\_pin\_record} \\
\texttt{TR-03} & \texttt{clean\_device\_verification\_report} \\
\hline
\end{tabular}
\end{center}

The release manifest binds the checkpoint and offline verifier.  The publication
record proves that exact package was made available through the approved versioned
channel.  The independent pin record must originate outside that mutable release
source.  The clean-device report binds the independently obtained pin and verified
private evidence-package digest without copying the private document.

\subsection{Distributed-network artefacts}

\begin{center}
\begin{tabular}{p{0.13\linewidth}p{0.75\linewidth}}
\hline
\textbf{Gate} & \textbf{Required artefact kinds} \\
\hline
\texttt{NET-01} & \texttt{validator\_response\_approval};
\texttt{validator\_removal\_exercise\_report} \\
\texttt{NET-02} & \texttt{independent\_node\_qualification\_report} \\
\texttt{NET-03} & \texttt{real\_network\_failure\_report} \\
\texttt{NET-04} & \texttt{multi\_device\_upgrade\_activation\_report} \\
\texttt{NET-05} & \texttt{public\_relay\_validation\_report};
\texttt{relay\_privacy\_observation} \\
\hline
\end{tabular}
\end{center}

The independent-node report must identify the exercised operational profiles.  The
real-network report must cover NAT, reconnect, delayed messages, relay outage and
partition recovery.  Upgrade evidence binds the exact signed release manifest,
checkpoint, governance approval and activation height.  Public-relay evidence
separately records availability/recovery and the observation that no Legal document
content crossed the relay.

\subsection{Security and recovery artefacts}

\begin{center}
\begin{tabular}{p{0.13\linewidth}p{0.75\linewidth}}
\hline
\textbf{Gate} & \textbf{Required artefact kinds} \\
\hline
\texttt{SEC-01} & \texttt{registry\_deployment\_record};
\texttt{live\_recovery\_revocation\_report} \\
\texttt{SEC-02} & \texttt{application\_security\_review};
\texttt{application\_security\_findings\_closure} \\
\texttt{SEC-03} & \texttt{consensus\_cryptography\_review};
\texttt{consensus\_findings\_closure} \\
\texttt{SEC-04} & \texttt{independent\_multi\_network\_soak\_report} \\
\hline
\end{tabular}
\end{center}

A review document alone is insufficient where findings exist; the corresponding
closure artefact is also mandatory.  Recovery evidence must show the independently
operated current-descriptor registry, epoch advancement, new-device authority,
lost-device revocation and stale-session rejection.  Soak evidence must represent
independent hosts and real production-like conditions rather than logical processes
on the development Mac.

\subsection{Qualified Legal-content artefacts}

\begin{center}
\begin{tabular}{p{0.13\linewidth}p{0.75\linewidth}}
\hline
\textbf{Gate} & \textbf{Required artefact kinds} \\
\hline
\texttt{LEGAL-01} & \texttt{qualified\_reviewer\_appointment} \\
\texttt{LEGAL-02} & \texttt{customer\_agreement\_qualified\_approval} \\
\texttt{LEGAL-03} & \texttt{supplier\_agreement\_qualified\_approval} \\
\texttt{LEGAL-04} & \texttt{contractor\_agreement\_qualified\_approval} \\
\texttt{LEGAL-05} & \texttt{employment\_document\_qualified\_approval} \\
\texttt{LEGAL-06} & \texttt{terms\_privacy\_dpa\_qualified\_approval} \\
\texttt{LEGAL-07} & \texttt{approved\_template\_library\_manifest} \\
\hline
\end{tabular}
\end{center}

Qualified approvals bind exact reviewed bytes, template family, supported
jurisdiction, reviewer identity and approval date.  The encrypted-library manifest
then proves that only those exact hashes were populated and records their current,
withdrawn or superseded state.  Generic AI-generated wording or an unqualified
internal review cannot produce these artefact classes.

\subsection{Controlled-release artefacts}

\begin{center}
\begin{tabular}{p{0.13\linewidth}p{0.75\linewidth}}
\hline
\textbf{Gate} & \textbf{Required artefact kinds} \\
\hline
\texttt{REL-01} & \texttt{assurance\_boundary\_publication\_record} \\
\texttt{REL-02} & \texttt{operator\_documentation\_publication\_record} \\
\texttt{REL-03} & \texttt{electronic\_signature\_claim\_publication} \\
\texttt{REL-04} & \texttt{internal\_multi\_device\_deployment\_report} \\
\texttt{REL-05} & \texttt{closed\_customer\_pilot\_report} \\
\texttt{REL-06} & \texttt{pilot\_metrics\_report} \\
\texttt{REL-07} & \texttt{pilot\_findings\_closure\_report} \\
\texttt{REL-08} & \texttt{frozen\_release\_manifest};
\texttt{independent\_release\_pin\_record} \\
\texttt{REL-09} & \texttt{final\_security\_approval};
\texttt{final\_legal\_product\_approval};
\texttt{final\_release\_approval} \\
\texttt{REL-10} & \texttt{preauthorization\_readiness\_report};
\texttt{production\_release\_authorization} \\
\hline
\end{tabular}
\end{center}

The publication records prove that the approved claim boundary and operating
documents reached the production channel.  Pilot completion and metrics are
separate because participation alone does not prove acceptable outcomes.  Final
approval requires three distinct artefacts and the configured distinct authority
roles.  Production authorization remains the last explicit decision after all
preceding evidence is available.  Its signed body must commit to the exact file
digest of a validated readiness report in which the other thirty-three gates pass
and only \texttt{REL-10} remains unmet; authorization is never inferred from a
completion percentage.

\subsection{Offline creation and countersigning}

The script
\texttt{backend/scripts/create\_glyphchain\_release\_evidence.py} supports two
signing operations and one assembly operation:

\begin{enumerate}
  \item \texttt{create} reads the exact release manifest and required retained
        artefact files, computes their SHA-256 digests, builds the canonical
        production evidence body and adds the first Ed25519 attestation;
  \item \texttt{attest} reads an existing entry, validates its unchanged body and
        appends a signature from a different authority key to a new output file.
  \item \texttt{inventory} validates one or more signed entries, rejects duplicate
        task identifiers and mixed release-manifest digests, sorts the entries by
        canonical task identifier and writes one versioned exact-release inventory.
\end{enumerate}

This allows delivery, security, Legal and release authorities to countersign on
separate machines without sharing private keys.  Each output uses exclusive
owner-only creation and refuses overwrite.  Signing keys, release manifests,
artefacts and existing entries must be regular files; symbolic links are refused.
Inputs have explicit size bounds, and the private key must be an owner-only
Ed25519 PEM file.  Encrypted PEM keys are supported by naming the environment
variable containing the password; the password is not placed in command-line
arguments, command history or evidence output.  The tool never writes the private
key, password or source artefact path into the evidence body.

The evaluator independently verifies all resulting signatures and authority roles.
Use of the creator does not appoint its signer as trusted and does not establish
independence: those properties remain in the separately administered trust-root
inventory.  The evaluator also compares the inventory's declared release digest
with the supplied manifest bytes before processing any non-empty inventory.
The trust-root inventory is itself governance-signed and exact-digest pinned; an
unsigned, altered or previously valid but no-longer-current authority set cannot
authorise evidence.

\subsection{Resource and schema boundaries}

The release evaluator reads a bounded release manifest, evidence inventory and
trust-root inventory.  It refuses symbolic links, non-regular files, invalid UTF-8,
malformed JSON, unsupported schemas and incorrectly typed top-level collections.
Evidence statements are limited to 4096 characters, signer identifiers and
artefact fields use bounded protocol syntax, and future or timezone-free
observation times fail closed.

The readiness report records the required kind list, evidence-body digest,
artefact-manifest digest, artefact count and attesting signer identifiers for every
task.  Auditors can therefore reproduce both the semantic contract and the exact
bytes used in the decision without exposing private Legal content.

\subsection{Verification and claim boundary}

Regressions prove that a wrong artefact kind cannot close a task, missing or
duplicate identifiers are rejected, symlink inputs cannot redirect either CLI,
and two separately held authority keys can create and countersign an entry that
passes the intended multi-role independent gate.  A complete synthetic inventory
continues to prove the algorithmic success path, while the real empty production
inventory remains blocked.

These mechanisms prove correct evidence handling.  They do not create or validate
the truth of an external event by themselves.  The responsible authority remains
accountable for the retained artefact, and independent or qualified status remains
subject to the configured governance and professional review boundary.
